\section{Governance and Trust Assumptions}
\label{sec:governance}

The protocol's trust surface is deliberately narrow. Market parameters are fixed at creation and identified by their hash; every settlement action requires a valid zero-knowledge proof registered in the facts registry (Section~\ref{sec:settlement}). Trust concentrates in the hub contract, which governs all capital flows, and in the liveness of off-chain infrastructure.

Beyond these trust surfaces, the protocol minimizes dependence on any single party. Markets are independent state machines (Section~\ref{sec:architecture}) that maintain their own state without shared global state; market creation from registered program hashes is permissionless (new program types are governance-gated), and so is participation as an operator, liquidator, or keeper. Intent cancellation, margin adjustment, and direct repayment are on-chain user actions requiring no operator cooperation; per-position enforcement proofs and escape settlement provide exit guarantees without operator participation~\cite{daloc-decent2026}.

Safety is structural: the system invariants of Section~\ref{sec:settlement} hold across all reachable states, and no state transition can occur without a valid zero-knowledge proof. Policy compliance is enforced at proof time rather than through execution-blocking (Section~\ref{sec:lp}). Liveness depends on each market's operator and on the chains, bridges, provers, and external venues used by that market; a configurable timeout bounds how long a market can stall before the operator is paused, per-position exit paths remain available, and after a longer escape timeout any pre-image holder can settle healthy, non-matured positions at oracle price. The decentralization properties, including correctness, completeness, and exit guarantees, are analyzed in the companion paper~\cite{daloc-decent2026}.

The architecture targets custody recoverability without reliance on the original off-chain operator. Users can submit cancellations, allowance revocations, share transfers, and recovery actions directly to the relevant on-chain contracts; for markets deployed on L2s, this may rely on that L2's L1 inbox, escape hatch, or forced-transaction mechanism. This does not mean that every economic action is immediately available from L1: liquidation, oracle updates, bridge movement, and external-venue unwinds still depend on the relevant chains and infrastructure being live. Each market defines its own data-availability and recovery policy, and the architecture admits progressive trust reduction as markets mature.

\subsection{Protocol Administration}

A multisig with timelock governs protocol contracts, which are upgradeable. Upgradeability permits response to critical vulnerabilities at the cost of trust in the multisig. Because core contracts control capital flows and settlement execution, an upgrade can alter routing and pause behavior across every market simultaneously. Pause and bridge-adapter whitelisting are also multisig actions; cross-chain routes are permissioned even though markets are not.

Proof programs (circuits and verification keys) and underwriter policies are immutable once registered; the upgrade mechanism cannot alter what conditions must be proven, only what effects the contract executes after a valid proof is verified. The residual trust implications of hub upgradeability are analyzed in the companion paper~\cite{daloc-decent2026}.

\subsection{Liveness Dependencies}

Settlement correctness is enforced by zero-knowledge proofs (Section~\ref{sec:settlement}), but liveness is not: if a market's off-chain infrastructure halts, no new settlement actions execute for that market. Existing positions remain safe (collateral stays locked, debt is unchanged, intent cancellation remains available), but undercollateralized positions go undetected by operator sweeps until the operator resumes. Either counterparty can independently detect and flag breached positions via per-position enforcement proofs at any time. To enable this, operators publish position pre-images to persistent off-chain storage at origination; the publication layer falls back from the operator's API to on-chain calldata to decentralized storage, ensuring counterparties retain the data needed for independent enforcement~\cite{daloc-decent2026}.

Each market defines a liveness timeout. Sweep proving scales linearly with portfolio size; a large book requires proportionally more proving time. If no operator sweep completes within that window, the operator is paused: no new originations until a sweep is delivered. Per-position operations by counterparties continue unaffected. The three-level exit hierarchy (per-position independence, operator pause, and escape timeout) is specified in the companion paper~\cite{daloc-decent2026}.

LTV limits, rates, and buffer sizes are produced off-chain by the risk engine. The serving computation is cryptographically verified via zero-knowledge proof (Section~\ref{sec:risk}), ensuring quotes derive correctly from the declared kernel and risk policy. The underwriter's committed policy (Section~\ref{sec:lp}) imposes outer bounds (rate floors, LTV ceilings, concentration caps) enforced at proof time. Neither mechanism validates the underlying model calibration or Monte Carlo assumptions between those bounds.

Cross-chain movement depends on the selected bridge infrastructure and state root relay
(Section~\ref{sec:crosschain}); the protocol inherits the bridge's and relay's liveness
and censorship properties. Authorization correctness is cryptographic
(ZK storage proofs over actual chain state), but proof-generation availability
depends on the data provider, and bridge execution depends on the source
and destination chains being live.
